\documentclass{article}
\usepackage[T1]{fontenc}
\usepackage{lmodern}
\usepackage{graphicx}
\usepackage{amsmath,amssymb}
\usepackage[a4paper,margin=2cm]{geometry}
\usepackage[absolute,overlay]{textpos}
\setlength{\TPHorizModule}{1mm}
\setlength{\TPVertModule}{1mm}
\usepackage[indonesian]{babel}
\usepackage{hyperref}
\setlength{\emergencystretch}{2em}
% unit: Halaman 31

\begin{document}

% === Halaman 31 (llm) ===

\section*{Serpent}

\begin{itemize}
    \item Salah satu finalis AES (Advanced Encryption Standard)
    \item Dirancang oleh Ross Anderson, Eli Biham, dan Lars Knudsen.
    \item Ukuran blok plainteks dan cipherteks di dalam Serpent adalah 128 bit
    \item Panjang kunci 128, 192, dan 256 bit
    \item Jumlah putaran 32 kali, setiap putaran melakukan operasi substitusi-permutasi
    \item Jumlah kotak S-box adalah 8 buah
\end{itemize}

\end{document}
